
Machine learning progress is increasingly measured through benchmark performance. Papers With Code (PWC) hosts thousands of benchmark leaderboards where research teams compete to achieve state-of-the-art (SOTA) results on shared evaluation datasets. Yet benchmarks themselves have lifecycles: a benchmark that defines the frontier today may be superseded when the community adopts a harder or more comprehensive successor. Understanding what predicts \textit{benchmark displacement} --- the event where one benchmark loses plurality status to another within the same task --- has implications for evaluation infrastructure design, research effort allocation, and interpretation of progress claims.

A natural candidate predictor is \textit{community breadth at benchmark introduction}: how many distinct research teams participated in a benchmark when it first achieved plurality status. Two competing mechanisms motivate this hypothesis. Under the \textbf{lock-in mechanism} (H1), broad early adoption creates a large network of stakeholders who have invested in the benchmark --- trained models, established baselines, published results --- generating switching costs that slow displacement (HR < 1). Under the \textbf{saturation mechanism} (H2), widespread adoption signals that the community has collectively treated the benchmark as solved, accelerating the search for harder successors (HR > 1). Both mechanisms are theoretically grounded: H1 follows Rogers' Diffusion of Innovations framework for opinion-leader lock-in \cite{rogers2003diffusion}; H2 aligns with framing in recent benchmarking discourse concerning benchmark overuse and community-driven replacement \cite{iclr2025benchmarking}.

The challenge in testing this hypothesis is measurement validity. Prior attempts using raw cumulative submission counts collapsed into time proxies (partial $r^2$ = 0.0011 with temporal controls), rendering Cox regression uninterpretable. A prior directional signal (HR = 0.871) appeared only for \textit{score velocity} ($\Delta score$\_lag1\_z) --- a fundamentally different construct from submission-count diversity --- at 22 displacement events, an underpowered sample where sampling noise dominates. The current paper tests \textit{submission-count diversity} at substantially greater statistical power (176 displacement events, 258 complete-case rows); the two results address different predictors and are complementary, not contradictory (see Section 5.4).

This paper makes three contributions:

\begin{enumerate}
\item \textbf{Methodological:} A 5-gate FAIL FAST pre-validation protocol (G0--G4) that separates data quality from hypothesis validity before any regression. Gates test join coverage (G0), time-independence of diversity predictors via partial $R^2$ (G1--G2), cross-benchmark variance (G3), and multicollinearity (G4). All five gates pass by wide margins, ruling out measurement failure as an explanation for the null.
\end{enumerate}

\begin{enumerate}
\item \textbf{Empirical:} A Cox regression test of community breadth \textit{as operationalized by paper submission count} as a displacement predictor, on 258 complete-case rows from the h-e2 panel (87 Koch et al. \cite{koch2021reduced} taxonomy tasks, 2015--2023, Papers With Code), with 176 displacement events, yielding HR = 1.006 (95\% CI = [0.846, 1.196], LRT p = 0.9495). This is a near-perfect null.
\end{enumerate}

\begin{enumerate}
\item \textbf{Scientific:} The combination of validated measurement and adequate statistical power means the null is scientifically informative. Submission-count community breadth diversity is ruled out as a predictor class for benchmark displacement timing in the Papers With Code ecosystem (2015--2023). Score velocity and institutional diversity remain untested at full power.
\end{enumerate}

